\subsection{Literal graphs and the nine variable tradeoffs}
The primary positive polynomial uses positive inputs
\[
 (b,g,h,s,w,\phi_+,i,j_+,n_+)\in\Pos^9,
\]
with $\phi=\phi_+-1$, $j=j_+-1$, $n=n_+-1$. The natural version uses
$(b_0,g_0,h_0,s_0,w_0,\phi,i_0,j,n)\in\N^9$ and sets
$b=b_0+1$, $g=g_0+1$, $h=h_0+1$, $s=s_0+1$, $w=w_0+1$, $i=i_0+1$.
The literal circuit orders these nine slots consistently. Both maps are
bijections of the declared domains, and every arithmetic shift is charged.

For each packing and majorant, the JSON records the source certificate,
source parameters, domain, ordered inputs, all fixed-leaf recipes, live DAG,
named register map, output index, operation counts, and exact-degree
certificate. A register name is not an additional unknown.
The fixed-binding file connects the recipe leaves to signed or unsigned
little-endian payloads; fixed power leaves are materialized too.

The balanced $W=170$, $p=169$, $K_c=4426$ certificate has
\[
 (r,n_0,c_0)=(18394,4776,341),
\]
so the quotient-cancelled degree formula gives $4\,400\,764$.
Removing its uniform target shift gives $K_c=4425$ and the odd threshold
branch; then $(r,n_0,c_0)=(18390,4775,341)$ and the degree is
$4\,399\,820$. Its lower degree comes at a higher gate cost because the
entire source no longer permits the same homogeneous sharing.

Two other explicitly retained tails use the same norm with different
$q$ and positive majorant $V_1$:
\begin{enumerate}
\item The simple tail has $q=64$ and
$V_1=k(8D^2-4D-1)GH$, proved above.
\item The inherited support tail has $q=2048^2$ and
$V_1=[T_6+256i(k^2+2)+2048]H$.
Its uniform majorant proof is retained as a pinned inert dependency.
\end{enumerate}
The quotient-cancelled tail is the one proved in detail in the preceding
section. For all three, exact cancellation-sensitive upper degrees are
matched by a nonzero top diagonal coefficient. The verifier expands the
small defect polynomial, not the huge final polynomial, and propagates
its exact degree through the remaining circuit identities. Random modular
replays are additional diagnostics, not the sole degree proof.

The finite source verifier proves exact equality of each evaluator with
its declared sparse source in the polynomial ring over four independent
port indeterminates. In particular, shared homogenization does not assume
the ports have special numerical values. Full independent circuit replays
then use the actual fixed residues and both domain translations.
The primary operation count is never replaced by the cost of an auxiliary
certificate graph or a denominator-multiplied polynomial.
